\section{Race-Condition Vulnerability}
As explained, race conditions can cause inconsistencies in data (incorrect or corrupt data), however they can also present themselves as a vulnerability to the program/system they are present within. For example, in a TOCTOU\footnote{Time-Of-Check, Time-Of-Use.} attack, the time between checking some data and using it differ (the race window), in the case of a race-condition being present, this may allow modification of the data to-be-fetched after the check, but before the time it is to be used. If this is an address to jump/return to, and is overwritten by way of this attack, it may be possible that control-flow could be hijacked, by setting another thread to concurrently the address (if locks or other synchronisation techniques are not present in the program checking and using the address) (\cite{geeksRCV}).

\subsection{Example}
\href{https://icode4.coffee/?p=1081}{I Code 4 Coffee: Hacking the Xbox 360 Hypervisor Part 2: The Bad Update Exploit} (\cite{badUpdate}) is a write-up of Ryan Miceli's (otherwise known as 'Grimdoomer') attempt at defeating the Xbox 360's hypervisor via software-only techniques.\\
The attack vector devised was vastly complex, containing many crucial components, without which this vector would be infeasible, however one crucial component was a race-condition vulnerability, allowing Grimdoomer to decompress an XKE payload ('Bootloader Update Payload') to a specified address in memory, stored in the \texttt{dec\_output\_buffer} pointer, wherein another thread was created 'hammering the \texttt{dec\_output\_pointer} to point into the last segment of the hypervisor', this alone would not allow for code execution in the hypervisor, however does allow overwriting, in this instance the last segment of the hypervisor, however this could instead be anywhere in the hypervisor to cause memory corruption in the hypervisor, with a modified XKE payload.